% 05_formato_json.tex - Manual Capítulo 5: Formato JSON de Escenarios
\chapter{Formato de Archivo de Escenario (JSON)}
\label{man:json}

\section{Estructura Básica}
\label{sec:json-estructura}

El formato es el del Apéndice A del enunciado. Claves en español (leídas solo en \code{scenario/});
el código interno usa nombres en inglés.

\screenshotfigure{json-escenario}{0.90}{Estructura del escenario JSON cargado y validado por la aplicación.}

\screenshotfigure{json-validacion}{0.85}{Mensaje de validación del escenario mostrando errores o confirmación correcta.}

\begin{minted}{json}
{
  "version": "1.0",
  "mapa": { "filas": 20, "columnas": 20 },
  "tipos_terreno": {
    "camino":  {"costo": 1, "transitable": true},
    "pasto":   {"costo": 2, "transitable": true},
    "bosque":  {"costo": 4, "transitable": true},
    "pantano": {"costo": 7, "transitable": true},
    "muro":    {"transitable": false}
  },
  "terreno": [
    ["camino", "camino", "pasto", "pasto"],
    ["camino", "muro",   "bosque", "pasto"],
    ...
  ],
  "bases": { "A": {"fila": 0, "columna": 0}, "B": {"fila": 19, "columna": 19} },
  "recurso": {"fila": 10, "columna": 10},
  "unidades": [
    {"id": "A1", "bando": "A", "tipo": "estandar", "fila": 5, "columna": 5},
    {"id": "B1", "bando": "B", "tipo": "estandar", "fila": 15, "columna": 15}
  ],
  "turno": "A",
  "juego": {
    "portador_recurso": null,
    "profundidad_maxima_minimax": 4,
    "limite_turnos": 100,
    "variante": "base",
    "turnos_recuperacion": 1,
    "inmunidad_robo": 1,
    "utilidad": {
      "entrega": 10,
      "recurso_cargado": 3,
      "distancia_objetivo": 1,
      "distancia_rival": 0.5
    }
  },
  "prueba": {
    "modo": "busqueda",
    "algoritmo": "A_STAR",
    "heuristica": "H2",
    "unidad_inicio": "A1",
    "objetivo": {"fila": 15, "columna": 15},
    "k": 4,
    "orden_acciones": "default"
  }
}
\end{minted}

\section{Campos Obligatorios}
\label{sec:campos-obligatorios}

\begin{longtable}{p{3cm}p{3cm}p{8cm}}
\toprule
\textbf{Campo} & \textbf{Tipo} & \textbf{Descripción} \\
\midrule
\endfirsthead
\toprule
\textbf{Campo} & \textbf{Tipo} & \textbf{Descripción} \\
\midrule
\endhead
\midrule
\multicolumn{3}{r}{\textit{Continúa}} \\
\endfoot
\bottomrule
\endlastfoot
\code{version} & string & Versión del formato (\code{"1.0"}) \\
\code{mapa.filas} & int $>0$ & Filas del grid \\
\code{mapa.columnas} & int $>0$ & Columnas del grid \\
\code{tipos\_terreno} & object & Catálogo: \code{nombre $\to$ \{costo, transitable\}} \\
\code{terreno} & string[][] & Matriz \code{filas $\times$ columnas} \\
\code{bases.A} & \code{\{fila, columna\}} & Posición base A \\
\code{bases.B} & \code{\{fila, columna\}} & Posición base B \\
\code{recurso} & \code{\{fila, columna\}} & Posición recurso (si nadie lo porta) \\
\code{unidades} & array & Lista: \code{id, bando, tipo, fila, columna} \\
\code{turno} & \code{"A"} | \code{"B"} & Bando que debe actuar \\
\code{juego.portador\_recurso} & string | \code{null} & id unidad que porta, o \code{null} \\
\bottomrule
\caption{Campos obligatorios del JSON}
\label{tab:json-obligatorios}
\end{longtable}

\section{Campos Opcionales del Equipo (RNF-02)}
\label{sec:campos-opcionales}

\begin{table}[H]
\centering
\caption{Campos opcionales del equipo (no alteran obligatorios)}
\begin{tabular}{lll}
\toprule
\textbf{Campo} & \textbf{Default} & \textbf{Descripción} \\
\midrule
\code{juego.limite\_turnos} & 100 & Límite de turnos; empate al alcanzar \\
\code{juego.profundidad\_maxima\_minimax} & 4 & Profundidad para Minimax/$\alpha$-$\beta$ \\
\code{juego.variante} & \code{base} & \code{base} o \code{dinamica} \\
\code{juego.turnos\_recuperacion} & 1 & Turnos propios de recuperación en variante dinámica \\
\code{juego.inmunidad\_robo} & 1 & Turnos del rival con inmunidad tras perder el recurso \\
\code{utilidad} & 10/3/1/0.5 & Pesos: \code{entrega}, \code{recurso\_cargado}, \code{distancia\_objetivo}, \code{distancia\_rival} \\
\code{prueba.algoritmo} & none & Algoritmo: \code{BFS}, \code{DFS}, \code{UCS}, \code{A\_STAR}, \code{BEAM}, \code{MINIMAX}, \code{ALPHA\_BETA} \\
\code{prueba.heuristica} & \code{H1} & \code{H1} (Manhattan) o \code{H2} (Obstacle-aware) \\
\code{prueba.k} & none & Ancho de haz (Beam) \\
\code{prueba.orden\_acciones} & \code{default} & \code{default} o \code{reversed} \\
\code{prueba.unidad\_inicio} & none & id unidad de inicio (búsqueda) \\
\code{prueba.objetivo} & none & \code{\{fila, columna\}} meta de búsqueda \\
\bottomrule
\end{tabular}
\end{table}
\label{tab:json-opcionales}

\section{Variante dinámica}
\label{sec:json-variante-dinamica}

La variante dinámica se activa con \code{"variante": "dinamica"} dentro de \code{juego} o
con el argumento CLI \code{--variant dinamica}. El costo del terreno determina los turnos
de recuperación de una unidad y el placaje puede robar el recurso respetando la inmunidad
configurada. Si no se especifica la variante, se usa \code{base}.


\section{Reglas de Validación (RF-34, RF-35)}
\label{sec:validacion}

\begin{enumerate}
\item \textbf{Atómica:} Se valida todo antes de construir; cualquier error $\to$
  \code{ScenarioError} con la lista completa; no se construye ningún escenario parcial.
\item \textbf{9 validaciones mínimas (RF-35) + 1 extra:}
  \begin{enumerate}
  \item Filas/columnas positivas y consistentes con matriz
  \item Terrenos usados existen en \code{tipos\_terreno}
  \item Transitables tienen costo positivo
  \item Bases, recurso, unidades dentro del mapa
  \item Ninguna unidad en celda no transitable
  \item IDs de unidad únicos
  \item Cada unidad en bando \code{A}/\code{a} o \code{B}/\code{b}
  \item Turno es \code{A}/\code{a} o \code{B}/\code{b}
  \item \code{portador\_recurso} es id existente o \code{null}
  \item \textbf{Extra:} Ninguna unidad comparte celda con otra
  \end{enumerate}
\item \textbf{Normalización:} Bandos \code{A}/\code{a} $\to$ \code{A}, \code{B}/\code{b} $\to$ \code{B};
  costos a \code{float}; terrenos no transitables sin \code{costo} tolerados.
\end{enumerate}

\section{Escenarios Incluidos}
\label{sec:escenarios-incluidos}

\begin{table}[H]
\centering
\caption{Escenarios en \code{scenarios/}}
\begin{tabular}{ll}
\toprule
\textbf{Archivo} & \textbf{Uso} \\
\midrule
\code{apendice\_a.json} & Ejemplo oficial 4$\times$4 \\
\code{caso\_01\_bfs\_vs\_dfs.json} & Caso 1 \\
\code{caso\_02\_bfs\_vs\_ucs.json} & Caso 2 \\
\code{caso\_03\_obstaculo\_heuristica.json} & Caso 3 \\
\code{caso\_04\_comparacion\_heuristicas.json} & Caso 4 \\
\code{caso\_05\_beam\_k.json} & Caso 5 \\
\code{caso\_07\_decision\_adversarial.json} & Casos 6--10 (5$\times$7) \\
\code{prueba\_15x15.json} & Rendimiento 15$\times$15 (3 unidades/bando) \\
\code{prueba\_20x20.json} & 20$\times$20 (3 unidades/bando) \\
\code{escenario\_03x03.json} .. \code{escenario\_30x30.json} & Escalabilidad \\
\bottomrule
\end{tabular}
\end{table}

\section{Crear un Escenario Nuevo}
\label{sec:nuevo-escenario}

\begin{enumerate}
\item Copia \code{apendice\_a.json} a \code{scenarios/mi\_escenario.json}
\item Modifica \code{mapa.filas}, \code{mapa.columnas}, \code{terreno}, \code{bases},
  \code{recurso}, \code{unidades}, \code{turno}
\item Ajusta \code{tipos\_terreno} y \code{juego} según necesites
\item Ejecuta: \code{uv run python -m tactical\_grid scenarios/mi\_escenario.json --verify}
\item Si hay errores, el mensaje lista TODOS los problemas; corrige y repite
\end{enumerate}

\textbf{No se requiere tocar código:} Cualquier JSON válido carga sin modificar el programa
(RF-36, RF-37, RNF-03, RNF-04).